% !PRINTED 322
\weberpage{322}

Eine Stammdiskriminante ist durch keine ungerade Quadratzahl teilbar. Ist aber $\Delta$ durch 4 teilbar, so muß
\weq{3}{\Delta \equiv 8,\ 12 \quad (\mathrm{mod}\ 16)}
sein. Denn wäre $\Delta \equiv 0, 4$ $(\mathrm{mod}\ 16)$, so würde nach Forthebung des Faktors 4 eine Diskriminante übrig bleiben.

Um also zu einem gegebenen $D$ den Stamm zu finden, hat man zunächst die größte ungerade Quadratzahl und dann noch eine so hohe Potenz von 4 abzusondern, daß $\Delta$ entweder ungerade und $\equiv 1$ (mod 4) oder $\equiv 8$, oder $\equiv 12$ (mod 16) wird.

4. \gesperrt{Eine eigentliche Diskriminante, die nicht in Faktoren zerlegbar ist, die selbst wieder Diskriminanten sind, heißt Primdiskriminante.}

Ist $p$ eine natürliche ungerade Primzahl, und wird das Zeichen $\pm$ so bestimmt, daß $\pm p \equiv 1$ (mod 4) ist, so gibt es folgende Primdiskriminanten:
\weq{4}{\pm p,\quad -4,\quad +8,\quad -8.}

5. \gesperrt{Jede Stammdiskriminante läßt sich auf eine und nur auf eine Weise in Primdiskriminanten zerlegen.}

Denn sondert man von $\Delta$ zuerst alle Faktoren $\pm p$ ab, so bleibt nur eine der Zahlen 1, $-4$, $+8$, $-8$ übrig.

\begin{center}
\textbf{§ 85. Das erweiterte \gesperrt{Legendre-Jacobische} Symbol.}
\end{center}

\gesperrt{Das Legendresche Symbol} $\left(\dfrac{m}{p}\right)$ hat, wenn $m$ eine beliebige positive oder negative von Null verschiedene Zahl, $p$ eine in $m$ nicht aufgehende natürliche Primzahl ist, den Wert $+1$, wenn $m$ quadratischer Rest, und den Wert $-1$, wenn $m$ quadratischer Nichtrest von $p$ ist (Bd. I, § 145).

Die Bedeutung des Symbols ist von \gesperrt{Jacobi} so erweitert worden, daß für $p$ auch eine zusammengesetzte Zahl gesetzt werden kann. Wenn aber $m$ eine Diskriminante ist, so empfiehlt sich bisweilen eine noch weitergehende Verallgemeinerung, die wir jetzt darlegen müssen\footnote{\gesperrt{Kronecker}, Berliner Sitzungsberichte, 30. Juli 1885. \gesperrt{Dirichlet-Dedekind}, Vorlesungen über Zahlentheorie, 4. Aufl., § 186. H.\ \gesperrt{Weber}, Göttinger Nachrichten, Januar 1893.}.
